\section{Sharp laws for general revelation profiles}\label{sec:profiles}
The convention at $\lambda(t)=0$ in each geometric sum is $\sum_{j=0}^{n-1}(1-\lambda(t))^j=n$. Every reference to all histories means legal visible histories and fresh private randomization, not future environment audit data.
The next result connects the general modulus to an all-strategy optimum; it is not used to prove Theorem~\ref{thm:modulus}. Let $a,\lambda:[0,1]\to[0,1]$ be continuous, with $a(0)=\lambda(0)=0$, $\lambda(t)>0$ for $t>0$, and $\lambda(1)=1$. A fixed hidden parameter is $\theta=st$, $s\in\{-1,1\}$. At each diagnostic call the observer selects a sign before the report, incurs the unobserved score $a(t)$ if wrong, and then receives the correct sign with probability $\lambda(t)$ and an inconclusive report otherwise. Reports are conditionally independent over calls. The score is never feedback. All model dependence resides in this positive raw kernel, not in the observer's program.

\begin{theorem}[A sharp task-visible acquisition profile]\label{thm:profile}
At every integer $n\ge1$, the minimax risk among all history-dependent strategies is
\begin{equation}\label{eq:profile-exact}
 D_n(a,\lambda)=\sup_{0\le t\le1}
       \frac{a(t)}{2n}\sum_{j=0}^{n-1}(1-\lambda(t))^j.
\end{equation}
One stationary two-label program attains this value at every horizon. Define
\begin{equation}\label{eq:profile-psi}
 \Psi_n(a,\lambda)=\sup_t a(t)
       \min\left\{1,\frac1{n\lambda(t)}\right\},
\end{equation}
with the minimum interpreted as one at $\lambda=0$. For the actual two-label marked family of this program,
\begin{equation}\label{eq:profile-modulus}
 K_n=\Psi_n(a,\lambda),\qquad
 \frac14\Psi_n\le D_n\le\frac12\Psi_n.
\end{equation}
The physical return time is one and $B_n=0$. Hence a continuum of raw experiments attains the general modulus against all histories, not only against a preselected controller.
\end{theorem}
\begin{proof}
At a fixed $t$, give the two signs equal prior probability. Until the first conclusive report, all visible histories and any private randomization have the same law for both signs, and the next sign choice has error probability $1/2$ under that prior. The probability of this event before call $j+1$ is $(1-\lambda(t))^j$. Later losses are nonnegative. This proves the lower bound for every history-dependent strategy and hence the minimax lower bound. A program initialized with a fair retained sign, unchanged on inconclusive reports and replaced on a conclusive report, has precisely that risk for each sign. It stores two labels and no clock, and so attains the lower bound simultaneously for all $n$ and $t$.

Relabel the two states as correct and wrong for mathematical evaluation under a fixed nonzero $t$. This relabeling is not information provided to the program. Its boundary arrays are
\[
 P_t=\begin{pmatrix}1&0\\\lambda(t)&1-\lambda(t)\end{pmatrix},
 \quad T_t=P_t,\quad r_t=(0,a(t))^T,\quad g_t=0.
\]
The same gain holds at zero since $a(0)=0$. A corrector with difference $v=u_{\rm wrong}-u_{\rm correct}$ has objective
$|a(t)-\lambda(t)v|+|v|/n$. Its minimum is
$a(t)\min(1,(n\lambda(t))^{-1})$: choose $v=0$ or $a(t)/\lambda(t)$. Conversely the vector
$z=(0,\min(1,(n\lambda(t))^{-1}))^T$
is feasible in~\eqref{eq:modulus-dual} and gives exactly that value. This proves the first identity in~\eqref{eq:profile-modulus}.

For $0<l\le1$, the geometric average is at most $\min(1,(nl)^{-1})$. It is at least half that quantity: $(1-l)^n\le(1+nl)^{-1}$ follows from $(1-l)^{-1}\ge1+l$ and the binomial inequality (with $l=1$ immediate), so
$[1-(1-l)^n]/(nl)\ge1/(1+nl)\ge\tfrac12\min(1,(nl)^{-1})$.
Multiply by $a(t)/2$, take the supremum and use the limiting convention at zero.
\end{proof}

Now introduce the same environment-owned fresh-mark and calibration task as in Theorem~\ref{thm:joint}, but replace its diagnostic amplitude $t$ and revelation probability $t^k$ by $a(t)$ and $\lambda(t)$. For clarity, a round has four charged calls: preparation, diagnostic decision and report, fresh-mark report, and terminal audit. The mark $V\sim\rho$ is independent and freshly prepared each round; its compact support $K$ is fixed. The audit scores
$\norm{V-c}^2+(b-\sigma\delta)^2$,
where $\sigma\in\{-1,1\}$ is fixed and never observed. A declared label has a fixed mark/calibration readout; the visible protocol phase is not a free remembered real number. Diagnostic and audit scores are not feedback. Risk in the next theorem is per round; division by four gives per physical call.

For $e\in[0,1/2]$ replace the diagnostic report probability by
$\lambda_e(t)=(\lambda(t)-e)_+$, with the same sign, action and unchanged audit wires. Define the visibility profile
\[
 A_*(u)=\sup\{a(t):\lambda(t)\le u\},\qquad 0\le u\le1.
\]
It tends to zero as $u\downarrow0$ by continuity and strict positivity of $\lambda$ away from zero. Suppose, when a simplified matching formula is desired, that
\begin{equation}\label{eq:profile-doubling}
 A_*(2u)\le C_a A_*(u)\quad(0<u\le1/2),\qquad
 q_{\lfloor M/2\rfloor}(\rho)\le C_q q_M(\rho)\quad(M\ge2).
\end{equation}
The exact sandwich below does not need these doubling conditions.

\begin{theorem}[A matched same-task multiresource law]\label{thm:profile-joint}
For every $n\ge1$, $M\ge2$, $0\le\delta\le1$, and $0\le e\le1/2$, the $M$-label minimax per-round risk satisfies
\begin{equation}\label{eq:profile-sandwich}
 q_M(\rho)+\delta^2+D_n(a,\lambda_e)
 \le\calR(4n,M,\delta,e)
 \le q_{\lfloor M/2\rfloor}(\rho)+\delta^2+D_n(a,\lambda_e).
\end{equation}
The source-to-target diagnostic deficiency $\varepsilon_*(e)$, with the virtual report due before any later action and without changing the environment audit, obeys
\begin{equation}\label{eq:profile-deficiency}
 e/2\le\varepsilon_*(e)\le e.
\end{equation}
Under~\eqref{eq:profile-doubling}, uniformly in the displayed budgets,
\begin{equation}\label{eq:profile-joint}
 \calR(4n,M,\delta,e)
 \asymp q_M(\rho)+\Psi_n(a,\lambda)+\delta^2+A_*(\varepsilon_*(e)).
\end{equation}
The upper program is parameter independent, uses $2\lfloor M/2\rfloor\le M$ labels, and is valid without an internally stored horizon. For a readout approximation of norm error at most $\xi$, its mark term is at most $(\sqrt{q_{\lfloor M/2\rfloor}(\rho)}+\xi)^2$. Program and workspace costs remain separate from this label count.
\end{theorem}
\begin{proof}
A fresh independent mark cannot be predicted from preceding rounds. Conditional on all earlier information, an $M$-state fixed-readout predictor has mark error at least $q_M(\rho)$; extra randomized readouts cannot improve the conditional squared error. Averaging the two signs of the invisible calibration shift gives a contribution at least $\delta^2$. At any fixed $t$, averaging the two diagnostic signs and applying the inconclusive-history argument of Theorem~\ref{thm:profile} gives $D_n$ before taking the supremum in $t$. These lower bounds are simultaneous under the product prior over the two independent signs, so they add. The loss is not a maximum of unrelated tasks.

For the upper bound retain the diagnostic sign together with a $J=\lfloor M/2\rfloor$-cell mark code. At the mark-report phase load the nearest-center index without changing the sign. At the audit output that center and zero for the calibration coordinate. Retain both labels through the next physical preparation. These rows give the upper risk in~\eqref{eq:profile-sandwich}; all phases are the visible four-call protocol. The numerical assertion follows from the $L^2$ triangle inequality.

The identity report map has error at most $e$ uniformly, giving the upper deficiency bound. By the intermediate value theorem choose $t$ with $\lambda(t)=e$. Both signs then have the same source report, whereas the two target report laws have total-variation distance $e$. Any common simulator output is at distance at least $e/2$ from one of them by the triangle inequality. It cannot consult the future audit. This proves~\eqref{eq:profile-deficiency}, not an unsupported exact equality with $e$.

Put $\Psi_{n,e}=\Psi_n(a,\lambda_e)$. Its definition gives
\[
 \Psi_{n,e}\ge\max\{\Psi_{n,0},A_*(e)\},\qquad
 \Psi_{n,e}\le\max\{2\Psi_{n,0},A_*(2e)\}.
\]
For the second inequality split $\lambda(t)>2e$, where $\lambda_e\ge\lambda/2$, from its complement. Thus $\Psi_{n,e}\asymp\Psi_{n,0}+A_*(e)$ under~\eqref{eq:profile-doubling}. The case $e=0$ follows directly. Equations~\eqref{eq:profile-deficiency} and~\eqref{eq:profile-doubling} compare $A_*(e)$ and $A_*(\varepsilon_*)$. The finite-horizon geometric-sum proof of Theorem~\ref{thm:profile} also applies to $\lambda_e$, even when it vanishes at positive $t$; its gain-zero and corrector identity are not asserted on those blind parameters. Use that finite-horizon estimate and quantization doubling to conclude.
\end{proof}

\begin{corollary}[Power and nonalgebraic rates]\label{cor:profile-examples}
For $a(t)=t^v$ and $\lambda(t)=t^k$, $v,k>0$,
\[
 \Psi_n=n^{-\min(v/k,1)},\qquad A_*(e)=e^{v/k}.
\]
For $\lambda(t)=t$ and $a(t)=1/\log(\mathrm e/t)$ on $t>0$, with $a(0)=0$,
\[
 \Psi_n=\frac1{\log(\mathrm e n)},\qquad
 A_*(e)=\frac1{\log(\mathrm e/e)}\quad(e>0).
\]
Both profiles satisfy the required doubling. Taking $\rho$ to be a Cantor product measure or a finite mixture with a positive highest-dimensional acquired component gives singular multiresource laws through its own $q_M(\rho)$.
\end{corollary}
\begin{proof}
For the powers split at $t=n^{-1/k}$. Above the split the expression is $t^{v-k}/n$, whose maximum is at the split if $v\le k$ and at one otherwise; below it $t^v$ is increasing. For the logarithmic example $a(t)$ increases and $a(t)/t=1/[t\log(\mathrm e/t)]$ decreases on $(0,1]$, so both pieces attain their maximum at $1/n$. The power doubling factor is $2^{v/k}$. In the logarithmic case it is at most $1+\log2$, since $e\le1/2$. The singular quantization assertions follow from the explicit prefix cover and retained submass proof of Theorem~\ref{thm:joint}; they do not turn a formal reachable set into an acquired distribution.
\end{proof}
The nonalgebraic example is deliberately outside the finite semialgebraic rate conclusion of Theorem~\ref{thm:algebraic}. Its exact two-label implementation and the converse in Theorem~\ref{thm:profile} remain valid. Accordingly a universal algebraic acquisition exponent cannot be inferred from compactness and task continuity alone.
